% title: Diagonal Quartic Modulo Prime
% short-title: Diagonal Quartic
% abstract: We prove that a four-variable diagonal quartic form has a non-trivial zero modulo every prime.
% subjclass: 11A07, 11T06
% keywords: prime modulus, fourth powers, cyclic groups, Lean 4
% author: Mario Hernández M.

\section{Proof of the theorem}

\begin{theorem}
Let \(p\) be a prime number. Then there exist integers \(x,y,z,t\), not all divisible by \(p\), such that
\[
p \mid x^4 - 2y^4 + 3z^4 + 4t^4.
\]
\end{theorem}

\begin{proof}
The formal statement is \lean{Biblioteca.Demonstrations.prime_dvd_diagonal_quartic_exists}.  We first prove the finite-field version
\lean{Biblioteca.Demonstrations.exists_zmod_diagonalQuarticForm_zero} for \(\mathbb F_p\), and then lift the representatives of
\(\mathbb Z/p\mathbb Z\) back to integers.

The cases \(p=2\) and \(p=3\) are immediate: respectively \((0,0,0,1)\) and \((0,0,1,0)\) are zeros.  Assume now
\(p>3\).  Work in the cyclic group \(\mathbb F_p^\times\) and quotient by the subgroup of fourth powers.  This quotient
has at most four elements, because the fourth-power map on a cyclic group has image of index
\(\gcd(4,p-1)\).

In this quotient write
\[
a=[2],\qquad r=[-1].
\]
Since \(r^2=1\), the following elementary intersection holds in every cyclic group of order at most four:
\[
\{1,a^2\}\cap \{a,ra^2,ra\}\ne\varnothing.
\]
Indeed, if \(a^2=1\), then either \(a=1\) or \(a\) is the unique non-trivial involution.  If \(a=1\), then
\(1=a\) belongs to both sets.  If \(a\ne 1\), then \(a\) is the unique non-trivial involution; since \(r^2=1\),
we have \(r=1\) or \(r=a\), and hence either \(1=ra^2\) or \(1=ra\).  Thus the two sets intersect.  If
\(a^2\ne 1\) and \(a\) has order \(3\), then the group has no non-trivial involution, so \(r=1\), and \(a^2\)
lies in both sets.  If \(a\) has order \(4\), then \(r=1\) or \(r=a^2\); in the first case \(a^2\) lies in both
sets, and in the second case \(ra^2=a^4=1\).  The Lean proof packages this reduced intersection as
\lean{Biblioteca.Demonstrations.cyclic_small_intersect}.

The two classes on the left are represented by the base values
\[
1=1^4+3\cdot0^4,\qquad
4=1^4+3\cdot1^4,
\]
while the three classes on the right are represented by
\[
2=2\cdot1^4-4\cdot0^4,\qquad
-4=2\cdot0^4-4\cdot1^4,\qquad
-2=2\cdot1^4-4\cdot1^4.
\]
If the two matching classes are \(L\) and \(R\), then \(R/L\) is a fourth power, say \(w^4\).  Multiplying the
corresponding left-hand representation by \(w^4\) gives
\[
(wx_0)^4+3(wz_0)^4=2y_0^4-4t_0^4,
\]
hence \( (wx_0)^4-2y_0^4+3(wz_0)^4+4t_0^4=0\) in \(\mathbb F_p\).  Because the right-hand base representation
always has \(y_0\) or \(t_0\) nonzero, the resulting quadruple is non-trivial.  This scaling step is formalized as
\lean{Biblioteca.Demonstrations.solutionOfClassEq}.
\end{proof}
